\documentclass{article}

\usepackage[a4paper,margin=1in]{geometry}

\usepackage{amsmath, amsthm, amssymb, amsfonts}
\usepackage{mathtools}
\usepackage{enumitem}
\usepackage{microtype}
\usepackage{xcolor}
\usepackage{fancyhdr}
\usepackage{graphicx}
\usepackage{hyperref}

% -------------------------------
% Fonts and Encoding
% -------------------------------
\usepackage{lmodern}
\usepackage[T1]{fontenc}
\usepackage[utf8]{inputenc}

% -------------------------------
% Theorem Environments
% -------------------------------
\theoremstyle{definition}
\newtheorem{definition}{Definition}[section]
\newtheorem{example}[definition]{Example}
\newtheorem{remark}[definition]{Remark}

\theoremstyle{plain}
\newtheorem{theorem}[definition]{Theorem}
\newtheorem{lemma}[definition]{Lemma}
\newtheorem{corollary}[definition]{Corollary}
\newtheorem{proposition}[definition]{Proposition}

% -------------------------------
% Header and Footer
% -------------------------------
\pagestyle{fancy}
\fancyhf{}
\rhead{\thepage}
\lhead{\textit{Mathematical Notes}}
\cfoot{}

% -------------------------------
% Custom Commands
% -------------------------------
\newcommand{\R}{\mathbb{R}}
\newcommand{\Z}{\mathbb{Z}}
\newcommand{\Q}{\mathbb{Q}}
\newcommand{\N}{\mathbb{N}}
\newcommand{\C}{\mathbb{C}}

\DeclareMathOperator{\Ker}{Ker}
\DeclareMathOperator{\Img}{Im}

% -------------------------------
% Examples & exercises
% -------------------------------
\newenvironment{solution}{\par\noindent\textbf{Solution.}\ }{\par}
\newenvironment{exercise}{\par\noindent\textbf{Exercise.}\ }{\par}

\begin{document}

\title{Examples and Exercises}
\author{Ramon}
\date{\today}

\maketitle

\section{Worked Examples}

\begin{example}[The field $\mathbb{R}$ and the vector space $\mathbb{R}^2$]
Let $\mathbb{K}=\mathbb{R}$ and let $V=\mathbb{R}^2$. Then $\mathbb{R}$ is a field under the usual addition and multiplication, and $V$ is a vector space over $\mathbb{R}$ under the usual vector addition and scalar multiplication
$$
\lambda(x,y)=(\lambda x,\lambda y).
$$
A basis of $V$ is given by
$$
e_1=(1,0),\qquad e_2=(0,1).
$$
Hence every vector $v\in V$ can be written uniquely in the form
$$
v=a_1e_1+a_2e_2,
\qquad a_1,a_2\in\mathbb{R}.
$$
For instance,
$$
(3,-2)=3e_1-2e_2.
$$
Define
$$
\langle (x_1,x_2),(y_1,y_2)\rangle=x_1y_1+x_2y_2.
$$
This is an inner product on $V$. Therefore two vectors are orthogonal precisely when their inner product is zero. For example,
$$
\langle (1,1),(1,-1)\rangle=1\cdot 1+1\cdot(-1)=0,
$$
so $(1,1)$ and $(1,-1)$ are orthogonal.
\end{example}

\begin{example}[A vector space over the complex field]
Let $\mathbb{K}=\mathbb{C}$ and let $V=\mathbb{C}^2$. Then $V$ is a vector space over $\mathbb{C}$. A basis is
$$
e_1=(1,0),\qquad e_2=(0,1).
$$
Thus every vector $z\in \mathbb{C}^2$ has the form
$$
z=a_1e_1+a_2e_2,
\qquad a_1,a_2\in\mathbb{C}.
$$
For example,
$$
(1+i,\,2-i)=(1+i)e_1+(2-i)e_2.
$$
This example emphasizes the distinction between the field and the vector space: the coefficients belong to $\mathbb{C}$, whereas the vectors belong to $\mathbb{C}^2$.
\end{example}

\begin{example}[A bilinear form that is symmetric but not an inner product]
Let $V=\mathbb{R}^2$ and define
$$
B((x_1,x_2),(y_1,y_2))=x_1y_1-x_2y_2.
$$
Then $B$ is bilinear and symmetric. However, $B$ is not an inner product, because it is not positive definite. Indeed,
$$
B((0,1),(0,1))=0\cdot 0-1\cdot 1=-1<0.
$$
Thus symmetry and bilinearity alone are not sufficient to produce an inner product.
\end{example}

\begin{example}[A function space with an inner product]
Let $V=P_2(\mathbb{R})$, the vector space of all real polynomials of degree at most $2$. Then $V$ is a vector space over $\mathbb{R}$. A basis is given by
$$
e_0(x)=1,\qquad e_1(x)=x,\qquad e_2(x)=x^2.
$$
Hence every polynomial $p\in V$ has a unique expression
$$
p(x)=a_0e_0(x)+a_1e_1(x)+a_2e_2(x)
=a_0+a_1x+a_2x^2,
\qquad a_0,a_1,a_2\in\mathbb{R}.
$$
Define
$$
\langle p,q\rangle=\int_{0}^{1} p(x)q(x)\,dx.
$$
This is an inner product on $V$. For example, let
$$
p(x)=1,\qquad q(x)=x-\frac12.
$$
Then
$$
\langle p,q\rangle=\int_0^1 \left(x-\frac12\right)\,dx
=\left[\frac{x^2}{2}-\frac{x}{2}\right]_0^1
=\frac12-\frac12=0.
$$
Therefore $p$ and $q$ are orthogonal in this inner-product space.
\end{example}

\begin{example}[Orthogonality in $\mathbb{R}^3$]
Let $V=\mathbb{R}^3$ with the usual inner product
$$
\langle (x_1,x_2,x_3),(y_1,y_2,y_3)\rangle=x_1y_1+x_2y_2+x_3y_3.
$$
Consider the vectors
$$
u=(1,2,1),\qquad v=(2,-1,0).
$$
Then
$$
\langle u,v\rangle=1\cdot 2+2\cdot(-1)+1\cdot 0=2-2+0=0.
$$
Hence $u$ and $v$ are orthogonal.
\end{example}

\section{Exercises with Solutions}

\begin{exercise}
Let $\mathbb{K}=\mathbb{R}$ and $V=\mathbb{R}^2$. Write the vector $(4,-3)$ as a linear combination of the standard basis vectors.
\end{exercise}

\begin{solution}
The standard basis of $\mathbb{R}^2$ is
$$
e_1=(1,0),\qquad e_2=(0,1).
$$
Hence
$$
(4,-3)=4e_1-3e_2.
$$
\end{solution}

\begin{exercise}
Explain why $\mathbb{R}$ and $\mathbb{R}^2$ are distinct sets, even though vectors in $\mathbb{R}^2$ use real numbers as coefficients.
\end{exercise}

\begin{solution}
The set $\mathbb{R}$ consists of scalars, whereas $\mathbb{R}^2$ consists of ordered pairs of scalars:
$$
\mathbb{R}^2=\{(x,y):x,y\in\mathbb{R}\}.
$$
Thus an element of $\mathbb{R}$ is a single real number, while an element of $\mathbb{R}^2$ is a pair of real numbers.
\end{solution}

\begin{exercise}
Let $V=\mathbb{R}^2$ and define
$$
B((x_1,x_2),(y_1,y_2))=2x_1y_1+3x_2y_2.
$$
Show that $B$ is bilinear.
\end{exercise}

\begin{solution}
Let $u=(x_1,x_2)$, $v=(x_1',x_2')$, $w=(y_1,y_2)$, and let $a,b\in\mathbb{R}$. Then
$$
au+bv=(ax_1+bx_1',\,ax_2+bx_2').
$$
Hence
\begin{align*}
B(au+bv,w)
&=2(ax_1+bx_1')y_1+3(ax_2+bx_2')y_2\\
&=a(2x_1y_1+3x_2y_2)+b(2x_1'y_1+3x_2'y_2)\\
&=aB(u,w)+bB(v,w).
\end{align*}
A similar computation shows
$$
B(w,au+bv)=aB(w,u)+bB(w,v).
$$
Therefore $B$ is bilinear.
\end{solution}

\begin{exercise}
Determine whether the bilinear form
$$
B((x_1,x_2),(y_1,y_2))=x_1y_2+x_2y_1
$$
on $\mathbb{R}^2$ is symmetric.
\end{exercise}

\begin{solution}
We compute
$$
B((x_1,x_2),(y_1,y_2))=x_1y_2+x_2y_1
$$
and
$$
B((y_1,y_2),(x_1,x_2))=y_1x_2+y_2x_1.
$$
Since multiplication in $\mathbb{R}$ is commutative,
$$
y_1x_2=x_2y_1,\qquad y_2x_1=x_1y_2.
$$
Thus
$$
B((y_1,y_2),(x_1,x_2))=x_2y_1+x_1y_2=B((x_1,x_2),(y_1,y_2)).
$$
Therefore $B$ is symmetric.
\end{solution}

\begin{exercise}
Show that the symmetric bilinear form
$$
B((x_1,x_2),(y_1,y_2))=x_1y_1-x_2y_2
$$
is not an inner product.
\end{exercise}

\begin{solution}
Take
$$
v=(0,1).
$$
Then
$$
B(v,v)=0\cdot 0-1\cdot 1=-1<0.
$$
Thus $B$ is not positive definite, hence it is not an inner product.
\end{solution}

\begin{exercise}
In $\mathbb{R}^2$ with the usual inner product, determine whether the vectors
$$
u=(2,1),\qquad v=(1,-2)
$$
are orthogonal.
\end{exercise}

\begin{solution}
The usual inner product is
$$
\langle (x_1,x_2),(y_1,y_2)\rangle=x_1y_1+x_2y_2.
$$
Hence
$$
\langle u,v\rangle=2\cdot 1+1\cdot(-2)=2-2=0.
$$
Since the inner product is zero, the vectors are orthogonal.
\end{solution}

\begin{exercise}
Let $V=P_1(\mathbb{R})$, the space of real polynomials of degree at most $1$, and define
$$
\langle p,q\rangle=\int_0^1 p(x)q(x)\,dx.
$$
Show that the polynomials
$$
p(x)=1,\qquad q(x)=x-\frac12
$$
are orthogonal.
\end{exercise}

\begin{solution}
We compute
$$
\langle p,q\rangle=\int_0^1 \left(x-\frac12\right)\,dx
=\left[\frac{x^2}{2}-\frac{x}{2}\right]_0^1
=\frac12-\frac12=0.
$$
Hence
$$
\langle p,q\rangle=0,
$$
so $p$ and $q$ are orthogonal.
\end{solution}

\begin{exercise}
Let $V=\mathbb{R}^3$. Express the vector
$$
v=(1,0,-2)
$$
as a linear combination of the standard basis vectors.
\end{exercise}

\begin{solution}
The standard basis of $\mathbb{R}^3$ is
$$
e_1=(1,0,0),\qquad e_2=(0,1,0),\qquad e_3=(0,0,1).
$$
Therefore
$$
(1,0,-2)=1e_1+0e_2-2e_3=e_1-2e_3.
$$
\end{solution}

\begin{exercise}
Let
$$
u=(1,1,1),\qquad v=(1,-1,0),\qquad w=(1,1,-2)
$$
in $\mathbb{R}^3$. Show that $u$ is orthogonal to both $v$ and $w$.
\end{exercise}

\begin{solution}
We compute
$$
\langle u,v\rangle=1\cdot 1+1\cdot(-1)+1\cdot 0=0,
$$
and
$$
\langle u,w\rangle=1\cdot 1+1\cdot 1+1\cdot(-2)=0.
$$
Hence $u$ is orthogonal to both $v$ and $w$.
\end{solution}

\begin{exercise}
Let $V=P_2(\mathbb{R})$ with basis
$$
e_0(x)=1,\qquad e_1(x)=x,\qquad e_2(x)=x^2.
$$
Write
$$
p(x)=2-3x+5x^2
$$
as a linear combination of the basis elements.
\end{exercise}

\begin{solution}
By direct comparison,
$$
p(x)=2e_0(x)-3e_1(x)+5e_2(x).
$$
\end{solution}

\end{document}
